\clearpage
\Article{348}
\textsc{Problem.} \emph{Given the equation}
\[
x^m+Ax^{m-1}+Bx^{m-2}+Cx^{m-3}+\text{etc.}+M=0
\]
\emph{whose roots are $x=a$, $x=b$, $x=c$, and so on, find an equation whose roots are $x=a^\tau$, $x=b^\tau$, $x=c^\tau$, and so on.}

\SourceLeaf{233}
\emph{First solution.} Find, by the known theorem, the sums of the roots of the given equation, of their squares, cubes, and so on, through the power with exponent $m\tau$. This also gives the sums of the roots of the required equation, and of their squares, and so on, namely $\Sigma a^\tau$, $\Sigma a^{2\tau}$, and so on. The coefficients can then be determined by the same theorem.

For practical calculation this solution is easier; but the following is better suited to our purpose, and also to showing that the coefficients of the required equation will be integers if the coefficients of the given equation are integers.

\Needspace{23\baselineskip}
\emph{Second solution.} Let $\theta$ be a primitive root of the equation $x^\tau=1$, and form the product of
\[
\begin{aligned}
&x^m+Ax^{m-1}+Bx^{m-2}+\text{etc.}\\
&x^m+A\theta x^{m-1}+B\theta\theta x^{m-2}+\text{etc.}\\
&x^m+A\theta\theta x^{m-1}+B\theta^4x^{m-2}+\text{etc.}\\
&\text{etc.}\\
&x^m+A\theta^{\tau-1}x^{m-1}+B\theta^{2\tau-2}x^{m-2}+\text{etc.}
\end{aligned}
\]
The roots of this product will therefore be
\[
\begin{array}{llll}
a,&\theta a,&\theta\theta a&\text{etc.}\\
b,&\theta b,&\theta\theta b&\text{etc.}\\
c,&\theta c,&\theta\theta c&\text{etc.}
\end{array}
\]
that is, the product will equal
\[
(x^\tau-a^\tau)(x^\tau-b^\tau)(x^\tau-c^\tau)\ldots
\]
and hence have the form
\[
x^{\tau m}+A'x^{\tau(m-1)}+B'x^{\tau(m-2)}+\text{etc.}
\]
\Needspace{14\baselineskip}
Now, replacing $x^\tau$ by $x$, we have
\[
x^m+A'x^{m-1}+B'x^{m-2}+\text{etc.}=(x-a^\tau)(x-b^\tau)(x-c^\tau)\ldots
\]
and hence
\[
x^m+A'x^{m-1}+B'x^{m-2}+\text{etc.}=0
\]
is the required equation. That $A',B'$, and so on, are not only rational numbers but also integers follows easily from the theory of the equation $x^\tau=1$.

Since we shall often use this operation below, by $(P,\rho^\tau)$
\SourceLeaf{234}
we shall denote the function which, on being set equal to zero, gives an equation whose roots are the powers with exponent $\tau$ of the roots of the equation $P=0$.

If $P\equiv Q$ for any modulus, then also $(P,\rho^\tau)\equiv(Q,\rho^\tau)$ for the same modulus.

\Needspace{10\baselineskip}
\Article{349}
\textsc{Theorem.} \emph{The coefficient of the term $x^n$ in $(P,\rho^\tau)$ is congruent, modulo $\tau$, to the coefficient of the term $x^{\tau n}$ in $P^\tau$, provided that $\tau$ is a prime number} (which, in this case, gives a third solution of the preceding problem).

\emph{Proof.} It follows from Chapter Six that every coefficient of the product
\[
(x^m+Ax^{m-1}+\text{etc.})(x^m+A\theta x^{m-1}+\text{etc.})\ldots
\]
has the following form after $\theta^\tau$ is replaced by unity:
\[
E+(1+\theta+\theta\theta+\text{etc.}+\theta^{\tau-1})F
\]
Now, if $\theta$ is regarded as a primitive root of the equation $x^\tau=1$, the whole product becomes $E$; but if $\theta=1$, the whole product becomes $P^\tau=E+\tau F$. Thus the coefficient of the term $x^{n\tau}$ in $P^\tau$ is congruent, modulo $\tau$, to the coefficient of the term $x^{n\tau}$ in $E$, that is, to the coefficient of the term $x^n$ in $(P,\rho^\tau)$.

\Needspace{10\baselineskip}
\Article{350}
\textsc{Theorem.} \emph{If $\tau$ is a prime number, then}
\[
(P,\rho^\tau)\equiv P\;(\mathit{mod.}\,\tau)
\]

\begingroup\postdisplaypenalty=10000\relax \emph{Proof.} Let the coefficient of the term $x^n$ in $(P,\rho^\tau)$ be $=N'$, and let the coefficient of the same term in $P$ be $=N$. Then, on putting
\[
P=x^m+Ax^{m-1}+\text{etc.}+Nx^n+\text{etc.}
\]
we have
\[
P^\tau\equiv x^{m\tau}+A^\tau x^{(m-1)\tau}+\text{etc.}+N^\tau x^{n\tau}+\text{etc.}\;(\mathrm{mod.}\,\tau)
\]
and hence, by the preceding article, $N'\equiv N^\tau\;(\mathrm{mod.}\,\tau)$. Thus, since $N^\tau\equiv N$, we obtain $N'\equiv N$. Q.\,E.\,D.\par\endgroup

It is also clear from this that $(P,\rho^\alpha)\equiv(P,\rho^{\alpha\tau})$ and $(P,\rho^\tau)\equiv(P,\rho^{\tau\tau})$, whence, in general,
\[
(P,\rho^\alpha)\equiv(P,\rho^{\alpha\tau^k})\;(\mathrm{mod.}\,\tau)
\]

\SourceLeaf{235}
\Needspace{10\baselineskip}
\Article{351}
\textsc{Theorem.} \emph{There is a value of the number $\nu$ less than $p^m$ such that the function $x^\nu-1$ can be divided by the given function $P$ of degree $m$, whose lowest term does not involve the indeterminate $x$, modulo $p$.}

\emph{Proof.} Divide by $P$ the sequence of functions $1,x,xx\ldots$ through $x^{p^m-1}$ as soon as their degree exceeds $m$. Since none can be divided by $P$ without a remainder, all the remainders can be put in the form
\[
Ax^{m-1}+Bx^{m-2}+\dotstwo+N
\]
with all coefficients positive and $<p$. But since they can never all be $=0$, there are clearly only $p^m-1$ functions, one of which each remainder must equal. Since, through the power of $x$ whose exponent is $p^m-1$, there are $p^m$ remainders, at least two must necessarily be the same. Let the same remainder result when $x^a$ and $x^{a+\nu}$ are divided by $P$, with $a+\nu<p^m$. Thus $x^{a+\nu}-x^a$ can be divided by $P$. Since, by hypothesis, $x$, and hence also $x^a$, is a function coprime to $P$, it follows that $x^\nu-1$ can also be divided by $P$. Q.\,E.\,D.

\emph{Corollary.} If $x^\nu-1$ can be divided by $P$, then $x^{k\nu}-1$ can also be divided by $P$, where $k$ denotes any integer.

\Needspace{10\baselineskip}
\Article{352}
\textsc{Theorem.} \emph{Keeping the notation of the preceding article, if $P$ is a prime function and $x^\nu$ is the lowest power which, after unity is subtracted, can be divided by $P$, then $\nu$ is either $=p^m-1$ or a divisor of this number, with the sole exception of the case $P=x$.}

\Needspace{26\baselineskip}
\emph{Proof.} Since $P$ is a prime function of degree $m$, there are $p^m-1$ different functions of degree less than $m$ (excluding the function $0$ from the total), all of which are coprime to $P$. Now, since $x^\nu$ is assumed to be the lowest power which leaves unity as remainder on division by $P$, it is clear that dividing all lower powers, from $1,x,\dotstwo$ through $x^{\nu-1}$, by $P$ gives $\nu$ different remainders, which we denote generally by $A$. If these exhaust all possible remainders, the theorem is proved; but if some are still absent from their number, let $B$ be any one of them. It is clear that the function $Bx^\nu$, divided by $P$, gives remainder $B$, and that, in general, $Bx^{\nu+k}\equiv Bx^k(\mathrm{mod.}\,P)$. But all the functions from $B$ through $Bx^{\nu-1}$
\SourceLeaf{236}
give remainders different from one another and from the remainders $A$. For if $Bx^\lambda\equiv Bx^{\lambda+\delta}(\mathrm{mod.}\,P)$, we would also have $1\equiv x^\delta(\mathrm{mod.}\,P)$, with $\delta<\nu$, contrary to the hypothesis; while if $Bx^\lambda\equiv x^\mu(\mathrm{mod.}\,P)$, we would have $B\equiv x^{\mu+\nu-\lambda}(\mathrm{mod.}\,P)$, and hence $B$ would be one of the remainders $A$, contrary to the hypothesis. We thus clearly have $\nu$ further new remainders. We may proceed in the same way (exactly as above, \textsection.\ .), and it will become apparent that the number $p^m-1$ of all possible remainders is either $=\nu$, or $=2\nu$, or $=3\nu$, or, in general, a multiple of the number $\nu$. Q. E. D.

\Needspace{10\baselineskip}
\Article{353}
It follows from the preceding theorem and the corollary to \textsection.\ 351 that every prime function of degree $n$ divides the function $x^{p^n-1}-1$ modulo $p$. Thus all functions of degree one, except the single one which is $\equiv x$, divide $x^{p-1}-1$, which is \textsc{Fermat's theorem}; all prime functions of degree two, that is, of the form $xx+Ax+B$, divide the function $x^{pp-1}-1$, and so on. Now let all divisors of the number $n$ be $n,\delta,\delta',\delta''$, and so on, .\ $1$. It is clear that $p^n-1$ is also divisible by $p^\delta-1$, $p^{\delta'}-1$, $p^{\delta''}-1$, and so on, $p-1$. Hence the function $x^{p^n-1}-1$ can be divided by \emph{all prime functions of degrees $n,\delta,\delta',\delta''$, and so on, down to the prime functions of degree one (excluding the function $x$)}, and therefore also, since all these functions are prime in themselves and hence coprime to one another, by \emph{the product of them all}. But this same product has degree $p^n-1$ (\textsection.\ 347), because the one function $x$ is absent. Thus it is clear that \emph{this product itself must be congruent to $x^{p^n-1}-1\;(\mathrm{mod.}\,p)$.}

\Needspace{20\baselineskip}
\Article{354}
\textsc{Theorem.} \emph{If the function $x^\nu-1$ is divisible by the function $P$, then}
\[
(P,\rho^{k\nu+t})\equiv(P,\rho^t)
\]
\emph{where $k,t$ denote any integers.}

\Needspace{29\baselineskip}
\emph{Proof.} Let
\[
P=x^m+Ax^{m-1}+Bx^{m-2}+\text{etc.}
\]
It is known that, if
\[
\frac{mx^{m-1}+(m-1)Ax^{m-2}+\text{etc.}}{x^m+Ax^{m-1}+\text{etc.}}
\]
is expanded in an infinite series of the form
\[
m\frac1x+\alpha\frac1{xx}+\beta\frac1{x^3}+\gamma\frac1{x^4}+\text{etc.}
\]
\SourceLeaf{237}
then $\alpha$ is the sum of the roots of the equation $P=0$, $\beta$ the sum of their squares, and so on. From this it follows without difficulty that the sums of the powers with exponents $\nu+1$, $\nu+2$, and so on, are congruent to the sum of the roots, of their squares, and so on. Hence, unless the modulus is equal to or less than the degree of the function $P$, we have
\[
(P,\rho^{\nu+1})\equiv P,\quad(P,\rho^{\nu+2})\equiv(P,\rho^2),\quad(P,\rho^{\nu+3})\equiv(P,\rho^3)\text{ etc.}
\]
We shall consider that case below.

\Needspace{20\baselineskip}
\Article{355}
\textsc{Theorem.} \emph{If, in the sequence}
\[
(P,\rho^0),\quad(P,\rho)\quad(P,\rho^2),\quad(P,\rho^3)\text{ etc.}
\]
\emph{the terms following the term in position $\nu$ are successively congruent to the first terms, then $x^\nu-1$ can be divided by $P$, provided that $P$ contains no factor more than once.}

\emph{Proof.} On putting $\frac{\mathrm{d}P}{\mathrm{d}x}=Q$, $Q$ will be a function coprime to $P$. Let
\[
\frac QP\equiv\frac Ax+\frac B{xx}+\frac C{x^3}+\text{etc.}
\]
Then, after the term $\frac N{x^\nu}$, there follows, by hypothesis,
\[
\frac A{x^{\nu+1}}+\frac B{x^{\nu+2}}+\frac C{x^{\nu+3}}+\text{etc.}
\]
Thus
\[
\frac QP\equiv\frac{Ax^{\nu-1}+Bx^{\nu-2}+\text{etc.}}{x^\nu-1}
\]
which shows that the function $x^\nu-1$ can be divided by $P$. Q. E. D.
